% The Julia chapter of the WFDB-x user's manual. It follows the sections of the Python
% chapter (chapters/python.tex), in the same order. Each \apientry of the section "The
% high-level API" is a function exported by the module WFDBX and defined in
% julia/src/highlevel.jl, julia/src/check.jl or julia/src/metadata.jl; its signature
% listing is copied from the method definition that takes the library state w::Wfdb first
% (each such function also has a method without w, with the same other arguments), and its
% params name every argument and keyword argument of the definition except w, in order.
% The examples of style "example" were run (each in a new module, after "using WFDBX", in
% a fresh copy of tests/data that is the current directory); check_manual.py does not run
% them.
\chapter{Julia}
\label{chap:julia}
\chapterlanguage{julia}

The Julia implementation of WFDB-x is the package \code{WFDBX} (module \code{WFDBX}), in
the directory \file{julia} of the repository: a library and the 33 command-line tools,
run through one driver script. It is a complete Julia implementation of the WFDB Software
Package: it uses only Julia and its standard libraries (no other package), and never calls
the C library.

It offers WFDB-x in three ways:
\begin{itemize}
\item \textbf{The high-level API} (\cref{sec:julia-api}): functions that read and write
  whole records, headers, annotation files and header metadata as Julia matrices, vectors
  and structures, in one call, with keyword arguments, and that throw a \code{WfdbError}
  (or, for some invalid arguments, an \code{ArgumentError}) on errors. They read signal
  files in large blocks, and are the fastest way to read a record. They are an addition
  of WFDB-x: the C package has no equivalent. They are the recommended interface for new
  code, analysis and notebooks.
\item \textbf{The WFDB library functions} under their C names (\code{isigopen},
  \code{getvec}, \code{getann}, \code{putann}, \ldots), with the behavior of the C library
  (\cref{sec:julia-library}). They are documented once, for all the languages, in
  \cref{chap:library}.
\item \textbf{The command-line tools} (\code{rdsamp}, \code{rdann}, \code{wrsamp},
  \code{xform} and 29 others), with the options, output and exit status of the C tools,
  run through the driver script \file{julia/bin/wfdbx.jl} (\cref{sec:julia-tools}). They
  are documented in \cref{chap:tools}.
\end{itemize}
The high-level functions return exactly the samples and annotations that the library
functions (and the C library) give.

% ==========================================================================================
\section{Installation and requirements}
\label{sec:julia-install}

\subsection{Requirements}

WFDB-x for Julia needs Julia 1.9 or later, on Linux, macOS or Windows (the Windows code
paths are not tested). It has no dependency beyond Julia's standard libraries: no WFDB
library and no libFLAC (the FLAC formats 508, 516 and 524 are read and written by
WFDB-x's own codec). Network access (\cref{sec:network}) uses the standard library
\code{Downloads}, which is loaded only when the first URL is read. The unit tests use the
standard libraries \code{Test} and \code{Sockets}, and the tests of \code{check\_record}
also run \cmd{python3}, with NumPy (\cref{sec:julia-testing}).

\subsection{Installing the package}

The package is installed from a copy of the repository, with the package manager:

\begin{lstlisting}[language=Julia]
using Pkg
Pkg.develop(path="/path/to/wfdb-x/julia")
\end{lstlisting}

In the package REPL, \code{] dev /path/to/wfdb-x/julia} does the same. The package is
\code{WFDBX}, version \code{1.2.0}, under the GNU General Public License, version 3 or
later, and is loaded with \code{using WFDBX}. The first \code{using WFDBX} precompiles
the package, running a workload of every part of the library and of every tool, so that
later sessions and tool runs load it without compiling anything. Changing a source file
of the package makes the next session precompile it again.

\subsection{The command-line tools}

The tools are run through the driver script \file{julia/bin/wfdbx.jl}: its first argument
names the tool, and the other arguments go to the tool. A small shell script, saved as
\file{wfdbx} in a directory of the \code{PATH}, makes the command shorter
(\cref{sec:julia-tools}):

\begin{lstlisting}[style=shell]
#!/bin/sh
# wfdbx: run a WFDB-x tool, as in `wfdbx rdsamp -r 100s -t 10`
exec julia --startup-file=no --project=/path/to/wfdb-x/julia \
    /path/to/wfdb-x/julia/bin/wfdbx.jl "$@"
\end{lstlisting}

\cmd{wfdbx -{}-version} then prints the version of WFDB-x, \code{WFDB-x 1.2.0}. The
docstrings of the package document each function (\code{?read\_record} in the REPL), with
examples.

\subsection{The sample record}
\label{sec:julia-resources}

The repository includes the sample record \code{100s}, the sample record of the WFDB
Software Package: the first minute of record 100 of the MIT-BIH Arrhythmia Database, two
ECG signals (\code{MLII} and \code{V5}) sampled at 360\,Hz, with the reference beat
annotations (annotator \code{atr}). Its files (\file{100s.hea}, \file{100s.dat},
\file{100s.atr}) are in the directory \file{tests/data} of a copy of the repository, with
the other test records of WFDB-x. The package does not install them: programs name that
directory as their input directory (\code{wfdbx\_set\_input\_dir("/path/to/wfdb-x/tests/data")},
or the keyword argument \code{input\_dir}). Work that uses the MIT-BIH Arrhythmia
Database cites it as described on PhysioNet.

% ==========================================================================================
\section{Quick start}
\label{sec:julia-quickstart}

This program reads the header, ten seconds of the signals and the annotations of the
sample record \code{100s}, from the directory \file{tests/data} of the repository (the
input directory of each call), and computes the mean heart rate. It then writes the first ten seconds
as a new record, \code{first10}, in the current directory, with a metadata entry in its
header.

\begin{lstlisting}[language=Julia,style=example]
using WFDBX

# The sample record 100s, in the directory tests/data of the repository.
const DATA = "tests/data"

# The header only: no samples are read.
hdr = read_header("100s"; input_dir=DATA)
println(hdr.nsig, " ", hdr.fs, " ", hdr.nsamples)     # 2 360.0 21600

# Seconds 10 to 20 of both signals, in ADC units (one column per signal).
rec = read_record("100s"; from="0:10", to="0:20", input_dir=DATA)
println(size(rec.samples))                          # (3600, 2)
x = physical(rec)                                   # in physical units
println(rec.samples[1, 1], " ", x[1, 1])            # 946 -0.39

# The beat annotations, and the mean heart rate.
ann = read_annotations("100s", "atr"; input_dir=DATA)
beats = ann.samples[[isqrs(c) != 0 for c in ann.codes]]
rr = (beats[end] - beats[1]) / ann.fs / (length(beats) - 1)
println(length(beats), " beats, ", round(60 / rr; digits=1), " beats per minute")
# 74 beats, 73.9 beats per minute

# Write the first 10 seconds as the record "first10" (first10.hea and first10.dat, in
# the current directory, the output directory), with the signal specifications of 100s.
r10 = read_record("100s"; to="0:10", input_dir=DATA)
write_record("first10", r10.samples, r10.fs; gains=200.0, baselines=1024, adczeros=1024,
             adcres=11, units="mV", names=signal_names(r10),
             metadata=["Source" => "MIT-BIH record 100"])

# Header metadata: "# Key: value" lines of the header (keys and values are bytes).
set_metadata("first10", "Age", "69")
for (key, value) in read_metadata("first10")
    println(String(key), " = ", String(value))
end
# Source = MIT-BIH record 100
# Age = 69
println(check_record("first10"))                    # no problems found
\end{lstlisting}

The program runs from the root of the repository. The input directory is given to each
call that reads \code{100s} (\code{input\_dir}), so the record \code{first10}, written to
the current directory, is then read from it under its plain name. The header
\file{100s.hea} gives no units: the program writes \code{"mV"} explicitly, since WFDB-x
never assumes units, a gain, an ADC resolution or a sampling frequency that a record does
not give (\cref{chap:behavior}).

% ==========================================================================================
\section{The high-level API}
\label{sec:julia-api}

The high-level API is made of the functions of this section, defined in the files
\file{highlevel.jl}, \file{check.jl} and \file{metadata.jl} of \file{julia/src} and
exported by the module \code{WFDBX}, of the types of their results
(\cref{sec:julia-results}), and of the exception type \code{WfdbError}
(\cref{sec:julia-warnings}). The docstrings of the package document each of them.

\subsection{Conventions}
\label{sec:julia-conventions}

\paragraph{The library state: \code{[w::Wfdb]}.} The state of the library (the open
record, the annotators, the input and output directories, network access, legacy mode) is
held by an object of the abstract type \code{Wfdb}; \code{Wfdb()} creates a new one, and
\code{lib()} returns the default object of the process. Each high-level function that uses
the state (all but \code{physical}, \code{signal\_names} and \code{lines}) has two
methods: one that takes a \code{Wfdb} object \code{w} as its first argument, such as
\code{read\_record(w, "100s")}, which uses the state of \code{w}, and one without it, with
the same other arguments, such as \code{read\_record("100s")}, which uses \code{lib()}.
The docstrings write this optional first argument \code{[w::Wfdb]}. The entries of this
section show the method that takes \code{w}, as it is defined; the method without it is
defined as \code{read\_record(record::ByteStr; kw...) = read\_record(lib(), record;
kw...)}. Separate \code{Wfdb} objects share no state (except what belongs to the process:
the count of write errors and the standard streams), so several records can be read at
once.

\begin{lstlisting}[language=Julia,style=example]
w = Wfdb()                              # a library state of its own
wfdbx_set_input_dir(w, "in")            # (the input directory of w only)
mkpath("in"); cp("100s.hea", "in/100s.hea"); cp("100s.dat", "in/100s.dat")
a = read_record(w, "100s")              # reads in/100s.hea, on w
b = read_record("100s")                 # reads ./100s.hea, on lib()
println(a.samples == b.samples)         # true
\end{lstlisting}

\paragraph{Keyword arguments.} The options of a function are keyword arguments, after a
semicolon, such as \code{read\_record("100s"; from="0:10", highres=true)}; each has a
default, given in the signature. An option whose default is \code{nothing} means ``not
given'': the default behavior or the current setting.

\paragraph{Record names and directories.} A record is named as in the other parts of
WFDB-x (\cref{chap:records,chap:files}). A full path, a \code{\textasciitilde} path or a
URL is used as given; any other name, such as \code{"100s"} or \code{"mitdb/100"}, is read
from the input directory and written to the output directory. Both are the current
directory unless they are set, with \code{wfdbx\_set\_input\_dir} and
\code{wfdbx\_set\_output\_dir} (\cref{sec:julia-settings}) or with the environment
variables \code{WFDBX\_INPUT\_DIR} and \code{WFDBX\_OUTPUT\_DIR}. So a record written under
a plain name is written to the output directory and read again from the input directory:
the two are the same unless one of them is set. URLs are read only when network access is
on. The readers take the keyword arguments \code{input\_dir} and \code{network}, and the
writers \code{dir}: they apply to that one call, and the previous settings are restored
when it returns, also when it throws. The signal files that a header names are read from
the header's directory.

A record name is given without the \file{.hea} extension (the readers, and
\code{write\_annotations}, also accept a trailing \code{.hea}; \code{write\_record} does
not). It is a full path, absolute (\code{"/data/mitdb/100"}) or in the home directory
(\code{"\textasciitilde/Documents/100s"}), used as given; a relative path
(\code{"100s"}, \code{"data/100s"}), read from the input directory if one is set,
otherwise from the current directory, and written to the output directory if one is set,
otherwise to the current directory (missing directories are created); or a URL
(\code{"https://physionet.org/files/mitdb/1.0.0/100"}), read only if network access is
on (a URL cannot be written). The \code{input\_dir} option applies to relative names
only, not to full paths or URLs. A writer given a name with a path creates the files in
that directory, and the header names the record and its signal file without the path
(\code{rec}, \code{rec.dat}), so that the record opens under any path that leads to it;
in legacy mode, as in the original, the header gives the name with its path. With the
\code{dir} option, the name must be a plain name: the files are created in \code{dir},
and the record opens from that directory under its plain name.

\paragraph{Times.} The keyword arguments \code{from} and \code{to} of \code{read\_record}
and \code{read\_annotations} are sample numbers (whole numbers, such as \code{3600} or
\code{3600.0}) or WFDB time strings, as \code{strtim} reads them (\code{"0:10"}, 10
seconds; \code{"1:02:03.5"}; \code{"s3600"}, sample 3600; \code{"e"}, the end of the
record; or a time of day in brackets, \code{"[14:00:00]"}). A time in seconds needs the
sampling frequency; when the record does not give it, such a time is an error.

\paragraph{Samples.} Samples are \code{Int32} values (C's \code{WFDB\_Sample}), in ADC
units, in a \code{Matrix\{Int32\}} with one row per sample and one column per signal:
\code{samples[i, s]} is sample \code{i} of signal \code{s} (both counted from 1).
Invalid samples are \code{WFDB\_INVALID\_SAMPLE} ($-32768$), and \code{NaN} in physical
units. Sample numbers and times are \code{Int64}, counted from 0 as in C. Signal numbers
are counted from 0 where the API gives them as numbers (\code{channels}, and the
\code{channels} field of a \code{Record}), as in C and the tools.

\paragraph{Per-signal and per-annotation values.} The per-signal keyword arguments of
\code{write\_record} (\code{gains}, \code{units}, \ldots) take one value for all the
signals, or a vector with one value per signal; the per-annotation fields of
\code{Annotations} (\code{symbols}, \code{chans}, \ldots) are empty or hold one value per
annotation. Any other length is an error, reported before anything is written.

\paragraph{Signal parameters.} No sampling frequency, gain, ADC resolution or units is
ever assumed. Where one is needed and the record does not give it, it must be given;
otherwise the function throws a \code{WfdbError}. A record whose header gives no sampling
frequency has \code{fs} 0, a signal without an ADC resolution has \code{adcres} 0, and a
signal without units has \code{units} \code{nothing}.

\paragraph{Text.} Text read from files (signal descriptions, units, info strings,
metadata, aux strings, mnemonics) is given as bytes, \code{Vector\{UInt8\}}, since WFDB
files are not UTF-8: every byte sequence round-trips unchanged. \code{String(x)} (or
\code{String.(v)} for a vector) converts bytes for printing; \code{String(x)} takes the
vector it converts, so \code{String(copy(x))} keeps \code{x}. Arguments accept
\code{String}s or bytes (the type \code{ByteStr}, \code{Union\{AbstractString,
AbstractVector\{UInt8\}\}}; \code{OptByteStr} also accepts \code{nothing}); a
\code{String} is taken as its bytes (UTF-8), as given. The names of records
(\code{Record.name}, \code{Header.name}, \ldots) and the time strings are \code{String}s.

\paragraph{Legacy mode.} With \code{legacy=true}, a function behaves as in the original
WFDB package, with its assumed signal parameters and its search for files
(\cref{chap:behavior}); the library's legacy mode is on for the duration of the call
only. With \code{false} (the default), the current setting of \code{wfdbx\_set\_legacy}
applies, which is off unless it was turned on.

\paragraph{State.} The functions that read and write records and annotation files call
\code{wfdbquit()} on the \code{Wfdb} object they use, before and after their work, also
when they throw, which closes anything opened with the WFDB library functions on that
object. The library's settings (legacy mode, getvec mode, time resolution of output
annotators, input and output directories, network access) are the same after a call as
before it: the high-level functions and \code{check\_record} save \code{getgvmode()} and
\code{getafreq()} when they start and restore them when they return, also when they
throw, since opening an EDF record adds \code{WFDB\_HIGHRES} to the getvec mode and sets
the time resolution to the record's sampling frequency (when none is set), as in C, and
\code{wfdbquit()} resets neither. The C-style functions keep the behavior of C. The metadata functions do not
call \code{wfdbquit()}: they open only the header, and leave the library's open files as
they are.

\paragraph{The examples} of this section run, each in a new module after \code{using
WFDBX}, in a copy of the test records of \file{tests/data} (among them \code{100s}),
which is the current directory and holds the records that the examples write; the input
and output directories are not set. The quick start (\cref{sec:julia-quickstart}) runs
instead from the root of the repository.

% ------------------------------------------------------------------------------------------
\subsection{Reading records}
\label{sec:julia-reading}

\apientry{read\_record}{Read the signals of a record}
\begin{lstlisting}[language=Julia,style=signature]
read_record(w::Wfdb, record::ByteStr; channels=nothing, from=0, to=nothing,
            highres::Bool=false, pad::Bool=false, input_dir::OptByteStr=nothing,
            network::Union{Nothing,Bool}=nothing, legacy::Bool=false,
            checksum_warning::Bool=true)
\end{lstlisting}
Reads the signals of a record from \code{from} up to \code{to} (exclusive), in ADC units,
together with the description of the record. The samples are those that \code{getvec}
returns (\cref{chap:library}), read in large blocks with \code{getvecs}. If the record
ends before \code{to}, the samples up to its end are returned. When the read reaches the
end of the record, the samples are compared with the checksums of the header
(\cref{sec:julia-warnings}). The getvec mode (\code{getgvmode()}) and the time resolution
of output annotators (\code{getafreq()}) are the same after the call as before it:
\code{highres} and \code{pad} apply to this call only.
\begin{params}
\item[record] Record name, without the \file{.hea} extension (a trailing \code{".hea"} is
  also accepted): a full path, absolute (\code{"/data/mitdb/100"}) or in the home
  directory (\code{"\textasciitilde/Documents/100s"}), used as given; a relative path
  (\code{"100s"}, \code{"data/100s"}), read from the input directory if one is set,
  otherwise from the current directory; or a URL
  (\code{"https://physionet.org/files/mitdb/1.0.0/100"}), read only if network access is
  on.
\item[channels] The signals to read, in the order given: a vector of signal numbers
  (\code{0} is the first signal) or names, such as \code{["V5"]} (as
  \code{signal\_names} gives them). \code{nothing}: every signal.
\item[from] The first sample: a sample number ($\ge 0$) or a time string. Default 0.
\item[to] The sample after the last one: a sample number or a time string.
  \code{nothing}: the end of the record.
\item[highres] \code{true} to read a multi-frequency record at its highest sampling
  frequency (\code{fs} of the result is then that frequency); \code{false}: at its frame
  frequency, as \code{getvec} does in low-resolution mode.
\item[pad] \code{true} to replace each invalid sample by the previous valid sample of
  the same signal (getvec mode \code{WFDB\_GVPAD}).
\item[input\_dir] The input directory for this call only (optional): a local directory,
  a \code{\textasciitilde} path or a URL, under which relative record names are read. It
  does not apply to full paths or URLs. \code{nothing}: the current setting.
\item[network] \code{true} to allow reading URLs for this call only, \code{false} to
  forbid it. \code{nothing}: the current setting.
\item[legacy] \code{true} to behave as the original WFDB package, for this call only.
\item[checksum\_warning] \code{true} (the default) to log a warning when the data do not
  match the checksums of the header; \code{false}: no warning.
\end{params}
\returns{A \code{Record} (\cref{tab:julia-record}), with the samples in ADC units.}
\raises{A \code{WfdbError} if the record cannot be read or has no signals (``can't read
record REC''), or its signals cannot be opened (``can't open signals of record REC''); if
a requested signal does not exist (``record REC has no signal C''); if \code{from} is
negative (``read\_record: from must not be negative''), if \code{from} or \code{to} is a
number that is not whole (``read\_record: from must be a whole number'') or beyond the
range of a 64-bit integer, or is neither a number nor a string (``read\_record: from must
be a sample number or a time string''); if \code{from} cannot be reached (``can't seek
to sample N of record REC''); if a time string needs the sampling frequency and the
header gives none (``sampling frequency of record REC is not specified''); or if a signal
file is shorter than its header says, or cannot be read or decoded (``can't read signal
data of record REC (getvec status S)''). Data that do not match the checksums of the
header are returned, with a warning (see \code{checksum\_warning}).}
\begin{lstlisting}[language=Julia,style=example]
# Seconds 10 to 20 of lead MLII, in physical units.
rec = read_record("100s"; channels=["MLII"], from="0:10", to="0:20")
mlii = physical(rec)[:, 1]
println(length(mlii), " ", rec.fs, " ", rec.sampfrom, " ", String(signal_names(rec)[1]))
# 3600 360.0 3600 MLII

# The first 3 samples of both signals, in ADC units, V5 first.
rec = read_record("100s"; channels=[1, 0], to=3)
println(rec.samples)                                # Int32[1011 995; 1011 995; 1011 995]
println(rec.channels, " ", rec.siginfo[1].gain, " ", rec.siginfo[1].baseline)
# [1, 0] 200.0 1024

# The whole record: the checksums are verified.
rec = read_record("100s")
println(size(rec.samples, 1), " ", rec.checksum_ok)  # 21600 true
\end{lstlisting}
A multi-frequency record (here the test record \code{mfreq}, whose signals have 1, 2
and 4 samples per frame at a frame frequency of 100\,Hz) is read at its frame frequency,
or at its highest sampling frequency:
\begin{lstlisting}[language=Julia,style=example]
rec = read_record("mfreq")
println(rec.fs, " ", size(rec.samples))             # 100.0 (600, 3)
rec = read_record("mfreq"; highres=true)
println(rec.fs, " ", size(rec.samples))             # 400.0 (2400, 3)
println([si.spf for si in rec.siginfo])             # Int32[1, 2, 4]
\end{lstlisting}

\apientry{read\_header}{Read the header of a record}
\begin{lstlisting}[language=Julia,style=signature]
read_header(w::Wfdb, record::ByteStr; input_dir::OptByteStr=nothing,
            network::Union{Nothing,Bool}=nothing, legacy::Bool=false)
\end{lstlisting}
Reads the header of a record, without opening its signal files. The result describes the
record and each of its signals, and holds the info strings and the metadata entries of the
header (from the \file{.hea} file only, as \code{read\_metadata} reads them, with its
warning for a key given more than once). For an EDF record, which has no \file{.hea}
file, and for a header read from the standard input (\code{"-"}), the metadata is empty,
without an error. For a multi-segment record, the master header is read, and
\code{nsamples} is the length of the whole record.
\begin{params}
\item[record] Record name, without the \file{.hea} extension (a trailing \code{".hea"} is
  also accepted): a full path, absolute (\code{"/data/mitdb/100"}) or in the home
  directory (\code{"\textasciitilde/Documents/100s"}), used as given; a relative path
  (\code{"100s"}, \code{"data/100s"}), read from the input directory if one is set,
  otherwise from the current directory; or a URL
  (\code{"https://physionet.org/files/mitdb/1.0.0/100"}), read only if network access is
  on.
\item[input\_dir] The input directory for this call only (optional), as for
  \code{read\_record}: relative record names are read under it. It does not apply to full
  paths or URLs.
\item[network] Network access for this call only, as for \code{read\_record}.
\item[legacy] \code{true} to read the header as the original WFDB package does, for this
  call only.
\end{params}
\returns{A \code{Header} (\cref{tab:julia-header}).}
\raises{A \code{WfdbError} if the header cannot be read (``can't read header for record
REC'').}
\begin{lstlisting}[language=Julia,style=example]
hdr = read_header("100s")
println(hdr.name, " ", hdr.fs, " ", hdr.nsig, " ", hdr.nsamples)   # 100s 360.0 2 21600
for si in hdr.siginfo
    println(String(si.desc), " ", String(si.fname), " ", si.fmt, " ", si.gain, " ",
            si.adcres, " ", si.units)
end
# MLII 100s.dat 212 200.0 11 nothing
# V5 100s.dat 212 200.0 11 nothing
println("[", String(hdr.comments[1]), "]")          # [ 69 M 1085 1629 x1]
println(hdr.base_time === nothing)                  # true: the header gives no start time
key, value = hdr.metadata[1]
println("[", String(key), "] [", String(value), "]")
# [Produced by xform from record 100, beginning at 0] [0]

# A copy of the header in the directory "in", read for this call only.
mkpath("in"); cp("100s.hea", "in/100s.hea")
println(read_header("100s"; input_dir="in").nsig)   # 2
\end{lstlisting}

\apientry{read\_annotations}{Read an annotation file}
\begin{lstlisting}[language=Julia,style=signature]
read_annotations(w::Wfdb, record::ByteStr, annotator::ByteStr; from=nothing,
                 to=nothing, fs::Union{Nothing,Real}=nothing,
                 input_dir::OptByteStr=nothing, network::Union{Nothing,Bool}=nothing,
                 legacy::Bool=false)
\end{lstlisting}
Reads the annotations of an annotation file, in file order, as parallel vectors. The
sampling frequency, needed to convert times in seconds, is that of the record's header,
or, for a record without a header or whose header gives none, the keyword argument
\code{fs}, as the option \opt{-F} of \code{rdann}.
\begin{params}
\item[record] Record name, without the \file{.hea} extension (a trailing \code{".hea"} is
  also accepted): a full path, absolute (\code{"/data/mitdb/100"}) or in the home
  directory (\code{"\textasciitilde/Documents/100s"}), used as given; a relative path
  (\code{"100s"}, \code{"data/100s"}), read from the input directory if one is set,
  otherwise from the current directory; or a URL
  (\code{"https://physionet.org/files/mitdb/1.0.0/100"}), read only if network access is
  on. Its header, if there is one, gives the sampling
  frequency.
\item[annotator] Annotator name: the suffix of the annotation file. \code{"atr"} reads the
  file \file{RECORD.atr}.
\item[from] The first time to read: a sample number or a time string. \code{nothing}: the
  beginning of the file.
\item[to] The time at which to stop: annotations at or after it are not read.
  \code{nothing}: the end of the file.
\item[fs] The sampling frequency, in Hz ($> 0$), to use when the record has no header or
  its header gives none. \code{nothing}: none.
\item[input\_dir] The input directory for this call only (optional), as for
  \code{read\_record}: relative record names are read under it. It does not apply to full
  paths or URLs.
\item[network] Network access for this call only, as for \code{read\_record}.
\item[legacy] \code{true} to behave as the original WFDB package (which assumes 250\,Hz
  for a record without a frequency), for this call only.
\end{params}
\returns{An \code{Annotations} (\cref{tab:julia-annotations}). Its \code{fs} is the
frequency in which the times are counted: the one written in the annotation file, if
there is one, or else the sampling frequency; it is 0 if none of them is known.}
\raises{A \code{WfdbError} if \code{fs} is given and is not $> 0$ (``read\_annotations:
fs must be > 0''); if the annotation file cannot be opened (``can't read annotator ANN for
record REC''); if a time string needs the sampling frequency and none is known
(``sampling frequency of record REC is not specified''); if \code{from} cannot be reached
(``can't seek in annotator ANN''); or if the file ends early or cannot be read before
\code{to} is reached (``can't read annotation file ANN of record REC (getann status S)'',
where S is the status that \code{getann} returned, such as $-3$ for an unexpected end of
file).}
\begin{lstlisting}[language=Julia,style=example]
ann = read_annotations("100s", "atr"; from="0:10", to="0:30")
println(length(ann), " ", ann.fs)                   # 24 360.0
println(ann.samples[1], " ", String(ann.symbols[1])) # 3862 N

# The whole file: the first annotation is a rhythm change, with an aux string.
ann = read_annotations("100s", "atr")
println(String(ann.symbols[1]), " ", ann.codes[1] == WFDBX.RHYTHM, " ", String(ann.aux[1]))
# + true (N

# A record without a header: fs gives the frequency of times in seconds (the library
# reports on the standard error that it cannot open the header).
write_annotations("beats", "qrs", Annotations(samples=[100, 460, 820],
                                              symbols=[b"N", b"N", b"N"]))
ann = read_annotations("beats", "qrs"; to="0:02", fs=360)
println(length(ann), " ", ann.fs)                   # 2 360.0
\end{lstlisting}

\apientry{physical}{Convert the samples of a record to physical units}
\begin{lstlisting}[language=Julia,style=signature]
physical(rec::Record; legacy::Bool=false)
\end{lstlisting}
Returns the samples of a \code{Record} in physical units, computed as $(\mathit{adc} -
\mathit{baseline}) / \mathit{gain}$ for each signal, with the gain and baseline of its
\code{siginfo}. Invalid samples become \code{NaN}. The result has the shape of
\code{rec.samples}. It uses no library state.
\begin{params}
\item[rec] The record, as \code{read\_record} returns it.
\item[legacy] \code{true} to behave as the original WFDB package, which converts the
  samples of an uncalibrated signal (gain 0) with a gain of 200. With \code{false}, the
  legacy mode of the default library object (\code{wfdbx\_set\_legacy}) decides.
\end{params}
\returns{A \code{Matrix\{Float64\}}, one row per sample and one column per signal.}
\raises{A \code{WfdbError} if the gain of a signal is 0 (``signal N of record REC is
uncalibrated'', where N is the number of the signal in the record, counted from 0),
except in legacy mode.}
\begin{lstlisting}[language=Julia,style=example]
rec = read_record("100s"; to=2)
println(physical(rec))          # [-0.145 -0.065; -0.145 -0.065]
\end{lstlisting}

\apientry{signal\_names}{Return the signal descriptions of a record}
\begin{lstlisting}[language=Julia,style=signature]
signal_names(rec::Record)
\end{lstlisting}
Returns the description of each signal of a \code{Record}, such as \code{"MLII"}, as
bytes, without the white space that a header may put at either end of it (empty for a
signal without a description). These are the names that \code{channels} of
\code{read\_record} accepts. It uses no library state.
\begin{params}
\item[rec] The record, as \code{read\_record} returns it.
\end{params}
\returns{A \code{Vector\{Vector\{UInt8\}\}}, one element per signal of \code{rec}.}
\raises{Nothing.}
\begin{lstlisting}[language=Julia,style=example]
rec = read_record("100s"; to=1)
println(String.(signal_names(rec)))     # ["MLII", "V5"]
\end{lstlisting}

% ------------------------------------------------------------------------------------------
\subsection{Writing records}
\label{sec:julia-writing}

\apientry{write\_record}{Write a header and a signal file}
\begin{lstlisting}[language=Julia,style=signature]
write_record(w::Wfdb, record::ByteStr, samples::AbstractMatrix{<:Integer}, fs::Real;
             fmt::Integer=16, gains=nothing, baselines=nothing, adczeros=0,
             adcres=nothing, units=nothing, names=nothing, samps_per_frame=nothing,
             base_time=nothing, comments=(), metadata=nothing, dir=nothing,
             legacy::Bool=false)
\end{lstlisting}
Writes samples in ADC units as a new record: the header \file{RECORD.hea} and one signal
file, \file{RECORD.dat}, which holds all the signals (so they share one format). Existing
files of these names are overwritten. The sampling frequency, the gain and ADC resolution
of each signal, and the units of each calibrated signal must be given: there are no
defaults. Every check except those of the writing itself comes before any file is
created, so an invalid argument writes nothing. The files are written as \code{putvec}
and \code{newheader} write them, followed by the info strings (\code{putinfo}).
\begin{params}
\item[record] Record name, without the \file{.hea} extension: a plain name
  (\code{"rec"}) or a relative path (\code{"data/rec"}), created in the output directory
  if one is set, otherwise in the current directory, missing directories included; or a
  full path, absolute or in the home directory (\code{"\textasciitilde/Documents/rec"}),
  used as given. With a path, the files are created in that directory, and the header
  names the record and its signal file without the path (\code{"rec"}, \code{"rec.dat"}),
  so that the record opens under any path that leads to it (in legacy mode, as in the
  original, the header gives the name with its path). A URL cannot be written. It must
  not be empty.
\item[samples] The samples in ADC units, one row per sample (per frame, with
  \code{samps\_per\_frame}) and one column per signal: integers within the range of
  \code{Int32}. \code{WFDB\_INVALID\_SAMPLE} is written as an invalid sample.
\item[fs] The sampling frequency in Hz: finite and $> 0$ (the frame frequency of a
  multi-frequency record).
\item[fmt] The storage format of all the signals, such as 16, 212, 24 or 516 (FLAC).
  Default 16.
\item[gains] The gain of each signal, in ADC units per physical unit (required): finite;
  0 marks an uncalibrated signal, and negative gains are allowed.
\item[baselines] The ADC value of 0 physical units of each signal. \code{nothing}: the
  signal's \code{adczero}.
\item[adczeros] The ADC value of 0\,V of each signal. Default 0.
\item[adcres] The ADC resolution of each signal, in bits, at least 1 (required).
\item[units] The physical units of each signal, such as \code{"mV"} (required for each
  calibrated signal; \code{nothing} for an uncalibrated one).
\item[names] The description of each signal. \code{nothing}: none.
\item[samps\_per\_frame] The number of samples per frame of each signal, integers $\ge 1$,
  for a multi-frequency record: each row of \code{samples} then holds
  \code{samps\_per\_frame[s]} consecutive columns for signal \code{s}. \code{nothing}: 1
  for every signal.
\item[base\_time] The start time and date, as \code{"hh:mm:ss dd/mm/yyyy"} (or a part of
  it, as \code{setbasetime} reads it). \code{nothing}: none.
\item[comments] The info strings to write at the end of the header, without their
  \code{\#}. Default: none.
\item[metadata] \code{Key: value} entries, written after the comments as
  \code{\# Key: value} (\code{\# Key:} for an empty value): a vector of \code{key =>
  value} pairs (or \code{(key, value)} tuples), written in that order, or a dictionary.
  Keys and values follow the rules of \code{set\_metadata}; no two keys may differ only in
  letter case, none may be the key of an entry in \code{comments}, and a value must not be
  \code{nothing}. \code{nothing}: none.
\item[dir] The directory (which must exist) in which to create the files, for this call
  only (the output directory of the call). \code{record} must then be a plain name,
  without a directory, and the record opens from that directory under its plain name.
  \code{nothing}: the output directory.
\item[legacy] \code{true} to behave as the original WFDB package, which assumes a gain of
  200, an ADC resolution based on the format and no units when they are not given, for
  this call only.
\end{params}
\returns{\code{nothing}.}
\raises{A \code{WfdbError}, before any file is created, if \code{record} is empty
(``write\_record: the record name must not be empty''); if there are no signals
(``write\_record: at least one signal is required''); if \code{gains}, \code{adcres} or
\code{units} is missing (``write\_record: gains must be specified'', ``write\_record: ADC
resolution(s) must be specified'', ``write\_record: units must be specified''); if a
\code{samps\_per\_frame} value is not an integer $\ge 1$ (``write\_record:
samps\_per\_frame must hold integers >= 1''), or the number of columns of \code{samples}
is not their sum (``write\_record: signals has N columns, but samps\_per\_frame adds up to
M''); if \code{fs} is not finite and $> 0$; if a gain is not finite (``write\_record:
gains must be finite''); if a vector of per-signal values has neither one value nor one
per signal (``write\_record: OPTION must have one value per signal''); if the number of
columns does not match the signals; if a sample lies outside the range of \code{Int32};
if \code{base\_time} is not a valid time (``write\_record: incorrect base time 'T'''); if
an entry of \code{metadata} is not valid (``write\_record: metadata key "a:b" must not
contain ':''', \ldots), or its key is also that of an entry in \code{comments}; or if
\code{dir} does not exist (``directory D does not exist'') or \code{record} contains a
directory while \code{dir} is given. A \code{WfdbError} also if a file cannot be created
(``can't create signal file for record REC''), or cannot be written completely, for
example on a full disk (``write error on record REC'').}
\begin{lstlisting}[language=Julia,style=example]
x = Int32[0 100; 10 110; 20 120; 30 130]            # 4 samples of 2 signals
write_record("two", x, 250; gains=[100, 200], adcres=12, units="mV",
             names=["a", "b"], base_time="12:00:00 01/01/2020",
             comments=["recorded at rest"], metadata=["Age" => "63"])
print(read("two.hea", String))
# two 2 250 4 12:00:00 01/01/2020
# two.dat 16 100/mV 12 0 0 60 0 a
# two.dat 16 200/mV 12 0 100 460 0 b
# #recorded at rest
# # Age: 63
println(read_record("two").samples == x)            # true

# A multi-frequency record: signal 1 has 2 samples per frame.
y = Int32[1 2 3; 4 5 6]                             # 2 frames: a, then b twice
write_record("mf", y, 100; gains=200, adcres=12, units="mV", samps_per_frame=[1, 2])
println(size(read_record("mf"; highres=true).samples))   # (4, 2)
\end{lstlisting}

\apientry{write\_record\_physical}{Write a record from samples in physical units}
\begin{lstlisting}[language=Julia,style=signature]
write_record_physical(w::Wfdb, record::ByteStr, samples::AbstractMatrix{<:Real}, fs::Real;
                      fmt::Integer=16, gains=nothing, baselines=nothing,
                      adczeros=0, adcres=nothing, units=nothing, names=nothing,
                      samps_per_frame=nothing, base_time=nothing, comments=(),
                      metadata=nothing, dir=nothing, legacy::Bool=false)
\end{lstlisting}
Writes a record as \code{write\_record} does, from samples in physical units. Each value
\code{v} is stored as $\mathrm{round}(v \cdot \mathit{gain}) + \mathit{baseline}$, with
halfway cases rounded away from zero; \code{NaN} is stored as an invalid sample. The gain
of every signal must be nonzero.
\begin{params}
\item[record] Record name, as for \code{write\_record}.
\item[samples] The samples in physical units, one row per sample (or frame) and one column
  per signal: real numbers; \code{NaN} for an invalid sample.
\item[fs] The sampling frequency in Hz, as for \code{write\_record}.
\item[fmt] The storage format, as for \code{write\_record}. Default 16.
\item[gains] The gain of each signal (required), finite and nonzero.
\item[baselines] The baseline of each signal, as for \code{write\_record}.
\item[adczeros] The ADC zero of each signal, as for \code{write\_record}.
\item[adcres] The ADC resolution of each signal (required), as for \code{write\_record}.
\item[units] The physical units of each signal (required), as for \code{write\_record}.
\item[names] The signal descriptions, as for \code{write\_record}.
\item[samps\_per\_frame] The samples per frame of each signal, as for
  \code{write\_record}.
\item[base\_time] The start time and date, as for \code{write\_record}.
\item[comments] The info strings, as for \code{write\_record}.
\item[metadata] The metadata entries, as for \code{write\_record}.
\item[dir] The directory in which to write the files, as for \code{write\_record}.
\item[legacy] \code{true} to behave as the original WFDB package (a gain of 0 is then
  taken as 200), for this call only.
\end{params}
\returns{\code{nothing}.}
\raises{A \code{WfdbError} in the cases of \code{write\_record}; if a gain is 0
(``write\_record: signal N is uncalibrated'', N counted from 0); or if a value is
infinite, or its ADC value lies outside the range of \code{Int32} (``write\_record:
signal N has a sample that is infinite or beyond the range of 32-bit ADC units'', N
counting the columns from 0). Nothing is written then.}
\begin{lstlisting}[language=Julia,style=example]
x = [0.0 1.0; 0.5 -1.0; NaN 0.25]                  # in mV; NaN: an invalid sample
write_record_physical("phys", x, 100; gains=200, adcres=12, units="mV", names=["a", "b"])
rec = read_record("phys")
println(rec.samples)            # Int32[0 200; 100 -200; -32768 50]
println(physical(rec))          # [0.0 1.0; 0.5 -1.0; NaN 0.25]
\end{lstlisting}

\apientry{write\_annotations}{Write an annotation file}
\begin{lstlisting}[language=Julia,style=signature]
write_annotations(w::Wfdb, record::ByteStr, annotator::ByteStr, ann::Annotations;
                  fs::Real=0.0, sort::Union{Nothing,Bool}=nothing, dir=nothing,
                  legacy::Bool=false)
\end{lstlisting}
Writes the annotations of an \code{Annotations} object to the annotation file
\file{RECORD.ANNOTATOR}, in the output directory or in \code{dir}, as \code{putann}
writes them. An existing file is overwritten. The type of each annotation is its mnemonic
in \code{ann.symbols} if that is not empty, and its code in \code{ann.codes} otherwise;
both may be given (as \code{read\_annotations} returns them) if they agree. As in C, the
subtype, channel and annotator number are written in one byte each (other values wrap
around), and auxiliary text longer than 255 bytes is truncated.
\begin{params}
\item[record] Record name, without the \file{.hea} extension (a trailing \code{".hea"}
  is also accepted): the record of the annotations (its header is not needed). A plain
  name or a relative path (\code{"data/100s"}) is created in the output directory if one
  is set, otherwise in the current directory; a full path, absolute or in the home
  directory (\code{"\textasciitilde/Documents/100s"}), is used as given. The file is
  \file{RECORD.ANNOTATOR} at that place (\file{data/100s.qrs}). A URL cannot be
  written.
\item[annotator] Annotator name: the suffix of the file, such as \code{"qrs"}.
\item[ann] The annotations (\cref{tab:julia-annotations}): \code{samples} (the times),
  and \code{symbols} or \code{codes}; \code{subtypes}, \code{chans}, \code{nums} and
  \code{aux} are optional. Each field is empty or has one element per annotation.
  \code{ann.record}, \code{ann.annotator} and \code{ann.fs} are not used.
\item[fs] 0 (the default), or a finite frequency $> 0$, which is written in the file as
  the frequency in which the annotation times are counted. With 0, readers use the
  record's sampling frequency. It applies to this file only: the time resolution of
  output annotators (\code{getafreq()}) is left as it was.
\item[sort] If the annotations are not in time order, the file is sorted when it is
  written: \code{true} always sorts, and \code{false} keeps the given order (without the
  library's advice to run \code{sortann}), for this call only. \code{nothing}: the
  environment variable \code{WFDBANNSORT} decides (sorting is on unless it is \code{0}).
\item[dir] The directory (which must exist) in which to create the file, for this call
  only. \code{record} must then be a plain name, without a directory: the file is
  \file{DIR/RECORD.ANNOTATOR}.
\item[legacy] \code{true} to behave as the original WFDB package, for this call only.
\end{params}
\returns{\code{nothing}.}
\raises{A \code{WfdbError}, before the file is created, if a field of \code{ann} is
neither empty nor has one element per annotation (``write\_annotations: FIELD must have
one entry per annotation (N), or none''); if \code{ann.symbols} and \code{ann.codes} are
both given and disagree (``write\_annotations: symbols and codes disagree at annotation
K'', K counted from 0); if \code{fs} is neither 0 nor a finite frequency $> 0$; or if
\code{dir} does not exist or \code{record} contains a directory while \code{dir} is given.
A \code{WfdbError} also if the file cannot be created (``can't create annotator ANN for
record REC'') or cannot be written completely (``write error on annotator ANN for record
REC'').}
\begin{lstlisting}[language=Julia,style=example]
ann = Annotations(samples=[100, 460, 820], symbols=[b"N", b"N", b"V"])
write_annotations("100s", "qrs", ann; fs=360)       # writes 100s.qrs
a = read_annotations("100s", "qrs")
println(a.samples, " ", String.(a.symbols))         # [100, 460, 820] ["N", "N", "V"]
println(a.codes == [WFDBX.NORMAL, WFDBX.NORMAL, WFDBX.PVC])   # true

# Out of order: sorted when the file is written (the library prints a message
# "Rearranging annotations for output annotator srt ...done!").
write_annotations("100s", "srt", Annotations(samples=[460, 100], codes=Int32[1, 1]); sort=true)
println(read_annotations("100s", "srt").samples)    # [100, 460]
\end{lstlisting}

% ------------------------------------------------------------------------------------------
\subsection{Checking records}
\label{sec:julia-checking}

\apientry{check\_record}{Check a record and its annotation files}
\begin{lstlisting}[language=Julia,style=signature]
check_record(w::Wfdb, record::ByteStr; annotators=String[],
             input_dir::OptByteStr=nothing, network::Union{Nothing,Bool}=nothing)
\end{lstlisting}
Checks the header of a record, its signal specifications, every sample of its signal
files, each segment of a multi-segment record, and the annotation files named, and
reports each problem found. A problem in the record is never an error of the function: it
is reported, with a severity (\code{"error"}: the record is damaged or does not follow
the WFDB format, and programs will read wrong data or fail; \code{"warning"}: the record
can be read, but something is missing or unusual, and may give wrong results;
\code{"info"}: a fact worth knowing, such as invalid samples). The problems, their codes
and their messages are the same in the seven languages (\cref{tab:python-issue-codes}).
The checks run in normal mode, even if legacy mode is on, and the library prints no
messages during the call (the previous settings are restored afterwards). They are
similar to those of Benjamin Moody's \code{wfdbcheck}.
\begin{params}
\item[record] Record name, without the \file{.hea} extension (a trailing \code{".hea"} is
  also accepted): a full path, absolute (\code{"/data/mitdb/100"}) or in the home
  directory (\code{"\textasciitilde/Documents/100s"}), used as given; a relative path
  (\code{"100s"}, \code{"data/100s"}), read from the input directory if one is set,
  otherwise from the current directory; or a URL
  (\code{"https://physionet.org/files/mitdb/1.0.0/100"}), read only if network access is
  on.
\item[annotators] The annotator names whose files are checked, such as \code{["atr"]}
  (one name may be given as a string). Default: none.
\item[input\_dir] The input directory for this call only (optional), as for
  \code{read\_record}: relative record names are read under it. It does not apply to full
  paths or URLs.
\item[network] Network access for this call only, as for \code{read\_record}.
\end{params}
\returns{A \code{CheckReport} (\cref{tab:julia-checkreport}): \code{report.ok} is
\code{true} if no problem is an error, \code{report.issues} lists the problems
(\code{Issue} objects, \cref{tab:julia-issue}), and \code{lines(report)} (or
\code{print(report)}) gives one line per problem, or \code{no problems found}.}
\raises{An \code{ArgumentError} if \code{record} is empty (``record must not be empty'')
or an annotator name is not a string (``annotator names must be strings''). A record that
cannot be read is reported (code \code{header\_unreadable}), not thrown as an error.}
\begin{lstlisting}[language=Julia,style=example]
report = check_record("100s"; annotators=["atr"])
println(report.ok)                                  # true: no errors
println(report)
# warning: signal 0: the units are not specified
# warning: signal 1: the units are not specified

# A damaged record: one sample of a 12-bit signal is overwritten with 2048.
write_record("damaged", reshape(Int32.(0:99), :, 1), 100; gains=100, adcres=12,
             units="mV", names=["ramp"])
dat = read("damaged.dat")
dat[21] = 0x00; dat[22] = 0x08      # sample 10 (2 bytes per sample in format 16): 2048
write("damaged.dat", dat)
report = check_record("damaged")
println(report.ok)                                  # false
for i in report.issues
    println(i.severity, " ", i.code, " ", i.where)
end
# error sample_range signal 0
# error checksum signal 0
println(lines(report)[1])
# error: signal 0: 1 sample is outside the ADC range (-2048 to 2047): sample 10 (value 2048)
\end{lstlisting}

\apientry{lines}{The text form of a check report}
\begin{lstlisting}[language=Julia,style=signature]
lines(r::CheckReport)
\end{lstlisting}
Returns the lines of a report of \code{check\_record}: one line per issue, in the order of
the issues (\code{"SEVERITY: WHERE: MESSAGE"}, without \code{"WHERE: "} for the record
itself), or the one line \code{"no problems found"}. \code{print(report)} and
\code{show(report)} write these lines, joined by line feeds. The lines are the same in the
seven languages.
\begin{params}
\item[r] The report, as \code{check\_record} returns it.
\end{params}
\returns{A \code{Vector\{String\}}, with at least one element.}
\raises{Nothing.}
\begin{lstlisting}[language=Julia,style=example]
for l in lines(check_record("100s"))
    println(l)
end
# warning: signal 0: the units are not specified
# warning: signal 1: the units are not specified
x = zeros(Int32, 10, 1)
write_record("clean", x, 100; gains=200, adcres=12, units="mV", names=["a"])
println(lines(check_record("clean")))               # ["no problems found"]
\end{lstlisting}

% ------------------------------------------------------------------------------------------
\subsection{Header metadata}
\label{sec:julia-metadata}

These functions read and change the metadata of a header: its info strings of the form
\code{\# Key: value}, such as \code{\# Age: 63} (\cref{chap:metadata}). In short: keys are
compared without regard to letter case (A-Z and a-z only); a key must not be empty and
must not contain \code{:}, \code{\#}, CR, LF or NUL, and a value must not contain CR, LF
or NUL; spaces, tabs, VT, FF and CR at either end of keys and values are removed; keys and
values are bytes (arguments \code{ByteStr}: a \code{String} is taken as its bytes;
results \code{Vector\{UInt8\}}). A change rewrites the header in one step, through a
temporary file in the directory of the destination, which then replaces it, keeping its
permissions and every other byte of it; the temporary file is removed on any failure. All
the arguments are checked first, so an error leaves the header unchanged, and nothing is
written if nothing changes (unless \code{dir} is given). Only the \file{.hea} file is read
(not a \file{.info} file); for a multi-segment record, it is the master header, unless a
segment is named. The header is opened as the library opens a header
(\code{wfdb\_open}), so in legacy mode it is searched for in the WFDB path. The entries
are returned as \code{MetadataEntry} pairs:
\begin{lstlisting}[language=Julia]
const MetadataEntry = Pair{Vector{UInt8},Vector{UInt8}}     # key => value
\end{lstlisting}
The functions take the keyword arguments \code{input\_dir} and \code{network} for reading
the header, and the functions that change it also \code{dir}, the directory in which to
write it. \code{dir} \code{nothing} or \code{""} changes the header that was read. These
functions do not call \code{wfdbquit()}. Invalid arguments throw an
\code{ArgumentError}, and a header that cannot be read or written a \code{WfdbError}
(\cref{sec:julia-warnings}). Metadata often holds information about patients: WFDB-x
stores it as plain text, exactly as given, and does not de-identify it.

\begin{phibox}
A header can hold identifying data about a patient (a name, a date of birth, a medical
record number). The metadata functions read and write such data as given. Records that
are shared must be de-identified first.
\end{phibox}

\apientry{read\_metadata}{Read the metadata of a record}
\begin{lstlisting}[language=Julia,style=signature]
read_metadata(w::Wfdb, record::ByteStr; input_dir::OptByteStr=nothing,
              network::Union{Nothing,Bool}=nothing)
\end{lstlisting}
Reads the metadata entries of the header of a record, in file order. An info string is an
entry if it contains a colon: the key is the text before the first colon, the value the
rest of the line, both without spaces at either end. Info strings without a colon, or
whose key is empty or contains \code{\#}, are free comments, which are not returned.
\begin{params}
\item[record] Record name, without the \file{.hea} extension (a trailing \code{".hea"} is
  also accepted): a full path, absolute (\code{"/data/mitdb/100"}) or in the home
  directory (\code{"\textasciitilde/Documents/100s"}), used as given; a relative path
  (\code{"100s"}, \code{"data/100s"}), read from the input directory if one is set,
  otherwise from the current directory; or a URL
  (\code{"https://physionet.org/files/mitdb/1.0.0/100"}), read only if network access is
  on.
\item[input\_dir] The input directory for this call only (optional), as for
  \code{read\_record}: relative record names are read under it. It does not apply to full
  paths or URLs.
\item[network] Network access for this call only, as for \code{read\_record}.
\end{params}
\returns{The entries, as \code{key => value} pairs of bytes, in file order. A key that
appears more than once (ignoring letter case) is returned once, with its first value, and
a warning is logged (\cref{sec:julia-warnings}).}
\raises{A \code{WfdbError} if the header cannot be found or read (``can't read header for
record REC''); an \code{ArgumentError} if \code{record} is empty (``record name must not be
empty'') or \code{"-"} (``the metadata of a header read from the standard input cannot be
read or changed'').}
\begin{lstlisting}[language=Julia,style=example]
write_record("pat", zeros(Int32, 360, 1), 360; gains=200, adcres=12, units="mV",
             names=["ECG"], comments=["recorded at rest"],
             metadata=["Age" => "63", "Sex" => "F"])
for (key, value) in read_metadata("pat")
    println(String(key), " = ", String(value))
end
# Age = 63
# Sex = F
for c in read_header("pat").comments
    println("[", String(c), "]")
end
# [recorded at rest]: the free comment is not an entry
# [ Age: 63]
# [ Sex: F]
mkpath("in"); cp("100s.hea", "in/100s.hea")
md = read_metadata("100s"; input_dir="in")          # reads in/100s.hea
println(length(md), " ", String(md[1].second))      # 1 0
\end{lstlisting}

\apientry{get\_metadata}{Return the value of a metadata entry}
\begin{lstlisting}[language=Julia,style=signature]
get_metadata(w::Wfdb, record::ByteStr, key::ByteStr; input_dir::OptByteStr=nothing,
             network::Union{Nothing,Bool}=nothing)
\end{lstlisting}
Returns the value of one metadata entry of a record. Letter case is ignored, and the
spaces at either end of \code{key} are removed: \code{"age"} finds \code{Age}. If the
header has more than one entry for \code{key}, the first is returned; as in
\code{read\_metadata}, a warning is logged for each key of the header that has more than
one entry.
\begin{params}
\item[record] Record name, as for \code{read\_metadata}.
\item[key] The key, such as \code{"Age"} (see \code{set\_metadata} for the valid keys).
\item[input\_dir] The input directory for this call only (optional), as for
  \code{read\_record}: relative record names are read under it. It does not apply to full
  paths or URLs.
\item[network] Network access for this call only, as for \code{read\_record}.
\end{params}
\returns{The value, as bytes (an empty vector for an entry with an empty value), or
\code{nothing} if the header has no entry for \code{key}.}
\raises{An \code{ArgumentError} if \code{key} is not a valid key (checked before the header
is read), or as \code{read\_metadata}.}
\begin{lstlisting}[language=Julia,style=example]
write_record("pat", zeros(Int32, 10, 1), 100; gains=200, adcres=12, units="mV",
             metadata=["Age" => "63", "Notes" => ""])
println(String(get_metadata("pat", "age")))         # 63
println(get_metadata("pat", "Notes"))               # UInt8[]
println(get_metadata("pat", "Weight") === nothing)  # true
\end{lstlisting}

\apientry{set\_metadata}{Set the value of a metadata entry}
\begin{lstlisting}[language=Julia,style=signature]
set_metadata(w::Wfdb, record::ByteStr, key::ByteStr, value::ByteStr;
             input_dir::OptByteStr=nothing, network::Union{Nothing,Bool}=nothing,
             dir::OptByteStr=nothing)
\end{lstlisting}
Sets the value of a metadata entry in the header of a record. If the header has an entry
for \code{key} (ignoring letter case), its line is replaced by \code{\# Key: value}
(\code{\# Key:} if the value is empty), at the same position and with the same line
ending, with the spelling of \code{key} given here. Otherwise the entry is added after the
last info string of the header (or, if it has none, after its last line that is not
blank), ending as the record line does. Every other byte of the header is kept.
\begin{params}
\item[record] Record name, as for \code{read\_metadata}. The header is changed where it
  was read, unless \code{dir} is given; a header read from a URL can only be written
  elsewhere, with \code{dir}.
\item[key] The key, such as \code{"Age"} or \code{"Date of birth"}. Spaces at either end
  are removed. It must not be empty, and must not contain \code{:}, \code{\#}, a line
  break or a NUL byte.
\item[value] The value, which may be a sentence. Spaces at either end are removed. It must
  not contain a line break or a NUL byte; it may be empty.
\item[input\_dir] The input directory for reading the header, for this call only.
\item[network] Network access for reading the header, for this call only.
\item[dir] An existing directory in which to write the changed header, under the same file
  name, instead of changing the header itself (the signal files are not copied).
  \code{nothing} or \code{""}: change the header that was read.
\end{params}
\returns{\code{nothing}.}
\raises{An \code{ArgumentError} if \code{record} is empty or \code{"-"}, if \code{key} or
\code{value} is not valid (``metadata key "a:b" must not contain ':''', ``value of metadata
key "Age" must not contain a line break'', \ldots), if the header has more than one entry
for \code{key} (``header of record REC has more than one entry for key KEY (lines N, M);
remove them first''), if it has no record line, or if it was read from a URL and
\code{dir} is not given. A \code{WfdbError} if the header cannot be read (``can't read
header for record REC''), if the changed header cannot be written (``can't write header
for record REC''), or if \code{dir} does not exist (``directory D does not exist'').}
\begin{lstlisting}[language=Julia,style=example]
write_record("pat", zeros(Int32, 10, 1), 100; gains=200, adcres=12, units="mV",
             metadata=["Age" => "63"])
set_metadata("pat", "age", "64")                    # replaces the line "# Age: 63"
set_metadata("pat", "Date of birth", "1960-01-01")  # adds a line
for l in eachline("pat.hea")
    startswith(l, "#") && println(l)
end
# # age: 64
# # Date of birth: 1960-01-01
mkpath("copy")
set_metadata("pat", "Age", "65"; dir="copy")        # writes copy/pat.hea
println(String(get_metadata("copy/pat", "Age")), " ", String(get_metadata("pat", "Age")))
# 65 64
\end{lstlisting}

\apientry{remove\_metadata}{Remove a metadata entry}
\begin{lstlisting}[language=Julia,style=signature]
remove_metadata(w::Wfdb, record::ByteStr, key::ByteStr; input_dir::OptByteStr=nothing,
                network::Union{Nothing,Bool}=nothing, dir::OptByteStr=nothing)
\end{lstlisting}
Removes a metadata entry from the header of a record. Every entry for \code{key} (ignoring
letter case) is removed, with its line ending, so this also repairs a header that has more
than one. Every other byte of the header is kept. Nothing is written if the header has no
entry for \code{key} (unless \code{dir} is given).
\begin{params}
\item[record] Record name, as for \code{set\_metadata}.
\item[key] The key, as for \code{set\_metadata}.
\item[input\_dir] The input directory for reading the header, for this call only.
\item[network] Network access for reading the header, for this call only.
\item[dir] The directory in which to write the changed header, as for
  \code{set\_metadata}.
\end{params}
\returns{The number of lines removed (0 if the header has no entry for \code{key}).}
\raises{As \code{set\_metadata} (a key that the header has more than once is not an error
here).}
\begin{lstlisting}[language=Julia,style=example]
write_record("pat", zeros(Int32, 10, 1), 100; gains=200, adcres=12, units="mV",
             metadata=["Age" => "63", "Sex" => "F"])
println(remove_metadata("pat", "SEX"))              # 1
println(remove_metadata("pat", "Weight"))           # 0: nothing is written
println(length(read_metadata("pat")))               # 1
\end{lstlisting}

\apientry{update\_metadata}{Make several changes to the metadata}
\begin{lstlisting}[language=Julia,style=signature]
update_metadata(w::Wfdb, record::ByteStr, changes; input_dir::OptByteStr=nothing,
                network::Union{Nothing,Bool}=nothing, dir::OptByteStr=nothing)
\end{lstlisting}
Makes several changes to the metadata of a record in one write. Each change is applied in
order, as \code{set\_metadata} does, or, if its value is \code{nothing}, as
\code{remove\_metadata} does. All the changes are checked first: if one is not valid, none
is made.
\begin{params}
\item[record] Record name, as for \code{set\_metadata}.
\item[changes] A vector (or tuple) of \code{key => value} pairs (or \code{(key, value)}
  tuples), or a dictionary; a value of \code{nothing} removes the entry. No two keys may
  differ only in letter case. The changes of a \code{Dict} are applied in its iteration
  order, which is not defined: a vector fixes the order in which new entries are added.
\item[input\_dir] The input directory for reading the header, for this call only.
\item[network] Network access for reading the header, for this call only.
\item[dir] The directory in which to write the changed header, as for
  \code{set\_metadata}.
\end{params}
\returns{\code{nothing}.}
\raises{As \code{set\_metadata}; an \code{ArgumentError} also if \code{changes} is not a
vector, a tuple or a dictionary of pairs (``metadata must be a dictionary, or a vector of
key => value pairs''), if a value is neither a string nor \code{nothing}, or if a key is
given more than once (``metadata key "AGE" is given more than once (keys are compared
without regard to letter case)'').}
\begin{lstlisting}[language=Julia,style=example]
write_record("pat", zeros(Int32, 10, 1), 100; gains=200, adcres=12, units="mV",
             metadata=["Age" => "63", "Dx" => "164889003"])
update_metadata("pat", ["Age" => "64", "Dx" => nothing, "Sex" => "M"])
println([String(k) => String(v) for (k, v) in read_metadata("pat")])
# ["Age" => "64", "Sex" => "M"]
\end{lstlisting}

% ------------------------------------------------------------------------------------------
\subsection{The result types}
\label{sec:julia-results}

The readers return objects of the types \code{Record}, \code{Header} and
\code{Annotations}, whose fields are listed in
\cref{tab:julia-record,tab:julia-header,tab:julia-annotations}; \code{check\_record}
returns a \code{CheckReport}, which holds \code{Issue} objects
(\cref{tab:julia-checkreport,tab:julia-issue}). The signals of a record are described by
\code{Siginfo} structures (\cref{tab:julia-siginfo}), which are those of the WFDB library
(\cref{sec:julia-library}). \code{Record}, \code{Header}, \code{CheckReport} and
\code{Issue} are immutable structures (the vectors and matrices they hold can be changed
in place); \code{Annotations} and \code{Siginfo} are mutable, with keyword constructors
whose defaults are empty vectors or C's zero values, such as \code{Annotations(samples=[100],
symbols=[b"N"])}. \code{length(ann)} is the number of annotations.

\begin{xltabular}{\linewidth}{@{}lL@{}}
\caption{The fields of \code{Record}, the result of \code{read\_record}}
\label{tab:julia-record} \\
\toprule
\tabhead Field & \tabhead Contents \\
\midrule
\endfirsthead
\toprule
\tabhead Field & \tabhead Contents \\
\midrule
\endhead
\bottomrule
\endfoot
\code{name::String} & the record name, as given to \code{read\_record} \\
\code{fs::Float64} & the sampling frequency of the returned samples, in Hz, or 0 if the
  header gives none; with \code{highres}, the highest sampling frequency of a
  multi-frequency record \\
\code{samples::Matrix\{Int32\}} & the samples in ADC units, one row per sample and one
  column per signal returned; invalid samples are \code{WFDB\_INVALID\_SAMPLE} \\
\code{siginfo::Vector\{Siginfo\}} & one per signal returned, in the order of the columns;
  \code{desc} has no white space at either end \\
\code{base\_time::Union\{Nothing,String\}} & the start time and date as text, to the
  second, such as \code{"[14:00:03 21/06/2024]"}; \code{nothing} if the header gives
  none \\
\code{base\_time\_ms::Union\{Nothing,Int64\}} & the start time in milliseconds after
  midnight (\code{"14:00:03.500"} gives \code{50403500}); \code{nothing} exactly when
  \code{base\_time} is \\
\code{base\_date::Union\{Nothing,NTuple\{3,Int\}\}} & the start date as \code{(year,
  month, day)}; \code{nothing} if the header gives no date \\
\code{sampfrom::Int64} & the number of the first sample returned \\
\code{comments::Vector\{Vector\{UInt8\}\}} & the info strings of the record: the comments
  of the header, and the lines of its \file{.info} file if it has one (see
  \code{getinfo}) \\
\code{channels::Vector\{Int\}} & the number, in the record, of each signal returned
  (\code{0} is the first signal) \\
\code{checksum\_ok::Bool} & \code{false} if the samples do not match the checksums of the
  header; checksums are checked only when a read reaches the end of the record, so a
  shorter read gives \code{true} \\
\end{xltabular}

\begin{xltabular}{\linewidth}{@{}lL@{}}
\caption{The fields of \code{Header}, the result of \code{read\_header}}
\label{tab:julia-header} \\
\toprule
\tabhead Field & \tabhead Contents \\
\midrule
\endfirsthead
\toprule
\tabhead Field & \tabhead Contents \\
\midrule
\endhead
\bottomrule
\endfoot
\code{name::String} & the record name, as given \\
\code{fs::Float64} & the sampling frequency in Hz, as \code{sampfreq} gives it; 0 if the
  header gives none \\
\code{counter\_freq::Float64} & the counter frequency (\code{getcfreq}) \\
\code{base\_counter::Float64} & the base counter value (\code{getbasecount}) \\
\code{nsamples::Int64} & the number of samples per signal, from the record line of the
  header (0 if the header does not give it) \\
\code{nsig::Int} & the number of signals \\
\code{siginfo::Vector\{Siginfo\}} & the signal specifications, one per signal; \code{desc}
  has no white space at either end \\
\code{base\_time}, \code{base\_time\_ms}, \code{base\_date} & the start time and date, as
  in \code{Record} \\
\code{comments::Vector\{Vector\{UInt8\}\}} & the info strings of the record, as in
  \code{Record} \\
\code{metadata::Vector\{MetadataEntry\}} & the metadata entries of the header, in file
  order, as \code{read\_metadata} returns them: from the \file{.hea} file only; empty,
  without an error, for an EDF record and for a header read from the standard input \\
\end{xltabular}

\begin{xltabular}{\linewidth}{@{}lL@{}}
\caption{The fields of \code{Annotations}, the result of \code{read\_annotations} and the
  argument of \code{write\_annotations}: parallel vectors, annotation \code{k} being
  \code{samples[k]}, \code{codes[k]}, \code{symbols[k]}, \ldots, in file order}
\label{tab:julia-annotations} \\
\toprule
\tabhead Field & \tabhead Contents \\
\midrule
\endfirsthead
\toprule
\tabhead Field & \tabhead Contents \\
\midrule
\endhead
\bottomrule
\endfoot
\code{record::String}, \code{annotator::String} & the record and annotator names that
  were read, such as \code{"100s"} and \code{"atr"}; not used by
  \code{write\_annotations} \\
\code{fs::Float64} & the frequency in which \code{samples} are counted, normally the
  sampling frequency of the record; 0 if unknown. \code{write\_annotations} does not use
  it \\
\code{samples::Vector\{Int64\}} & the annotation times, in samples \\
\code{codes::Vector\{Int32\}} & the annotation codes, such as \code{WFDBX.NORMAL} or
  \code{WFDBX.PVC} \\
\code{symbols::Vector\{Vector\{UInt8\}\}} & the mnemonics of the codes, as \code{annstr}
  returns them, such as \code{b"N"}, \code{b"V"} or \code{b"+"}; when writing, an empty
  mnemonic means ``use \code{codes}'' \\
\code{subtypes}, \code{chans}, \code{nums} (\code{Vector\{Int32\}}) & the subtype,
  channel number and annotator number of each annotation \\
\code{aux::Vector\{Union\{Nothing,Vector\{UInt8\}\}\}} & the auxiliary text of each
  annotation, without its length byte, or \code{nothing} where there is none \\
\end{xltabular}

\begin{table}[htbp]
\centering
\caption{The fields of \code{CheckReport}, the result of \code{check\_record}}
\label{tab:julia-checkreport}
\begin{tabularx}{\linewidth}{@{}lL@{}}
\toprule
\tabhead Member & \tabhead Contents \\
\midrule
\code{record::String} & the record name, as given \\
\code{ok::Bool} & \code{true} if no issue has the severity \code{"error"} \\
\code{issues::Vector\{Issue\}} & the problems, in the order of the checks \\
\code{lines(report)} & the text form of the report: one line per issue, or the one line
  \code{no problems found} \\
\code{print}, \code{show} & write the lines of \code{lines(report)}, joined by line feeds
  (no final line feed) \\
\bottomrule
\end{tabularx}
\end{table}

\begin{table}[htbp]
\centering
\caption{The fields of \code{Issue}, a problem found by \code{check\_record}}
\label{tab:julia-issue}
\begin{tabularx}{\linewidth}{@{}lL@{}}
\toprule
\tabhead Member & \tabhead Contents \\
\midrule
\code{severity::String} & \code{"error"}, \code{"warning"} or \code{"info"} \\
\code{code::String} & a stable identifier of the kind of problem, such as
  \code{"checksum"} or \code{"gain\_missing"} (\cref{tab:python-issue-codes}) \\
\code{where::String} & where the problem is: \code{""} (the record), \code{"signal N"},
  \code{"segment SEG"}, \code{"segment SEG, signal N"} or \code{"annotator ANN"} \\
\code{message::String} & the description of the problem \\
\code{string(issue)} & \code{SEVERITY: WHERE: MESSAGE}, or \code{SEVERITY: MESSAGE} when
  \code{where} is empty \\
\bottomrule
\end{tabularx}
\end{table}

\begin{xltabular}{\linewidth}{@{}lL@{}}
\caption{The fields of \code{Siginfo} (C's \code{WFDB\_Siginfo}), the specification of a
  signal}
\label{tab:julia-siginfo} \\
\toprule
\tabhead Field & \tabhead Contents \\
\midrule
\endfirsthead
\toprule
\tabhead Field & \tabhead Contents \\
\midrule
\endhead
\bottomrule
\endfoot
\code{fname::Union\{Nothing,Vector\{UInt8\}\}} & the name of the signal file, as the
  header gives it \\
\code{desc::Union\{Nothing,Vector\{UInt8\}\}} & the signal description, such as
  \code{b"MLII"} \\
\code{units::Union\{Nothing,Vector\{UInt8\}\}} & the physical units, such as \code{b"mV"};
  \code{nothing} if not specified \\
\code{gain::Float64} & the gain, in ADC units per physical unit (0: uncalibrated) \\
\code{initval::Int32} & the value of the first sample \\
\code{group::UInt32} & the signal group number (the signals of one group share one
  file) \\
\code{fmt::Int32} & the storage format, such as \code{16}, \code{212} or \code{516} \\
\code{spf::Int32} & the number of samples per frame (more than 1 for the faster signals of
  a multi-frequency record) \\
\code{bsize::Int32} & the block size, in bytes (0 for ordinary files) \\
\code{adcres::Int32} & the ADC resolution, in bits (0: not specified) \\
\code{adczero::Int32} & the ADC output for an input of 0\,V \\
\code{baseline::Int32} & the ADC value that corresponds to 0 physical units \\
\code{nsamp::Int64} & the number of samples (0: not specified) \\
\code{cksum::Int32} & the 16-bit checksum of the samples \\
\end{xltabular}

% ------------------------------------------------------------------------------------------
\subsection{Warnings and errors}
\label{sec:julia-warnings}

\paragraph{Exceptions.} The high-level functions report errors by exceptions
(\cref{tab:julia-exceptions}). \code{WfdbError} is the exception type of WFDB-x; its field
\code{msg} (a \code{String}) describes the failed operation, such as ``can't read record
REC'', and \code{showerror} prints \code{WfdbError: MESSAGE}. The metadata functions and
\code{check\_record} throw Julia's \code{ArgumentError} (whose field is \code{msg} too)
for an invalid argument, and a \code{MethodError} is thrown, as for any Julia function,
for an argument of a type that no method accepts (such as a number as the value of
\code{set\_metadata}). Each function checks its arguments before it creates or changes a
file, so an invalid argument leaves the files as they were. The library may also print a
message of its own on the standard error when it cannot open a file (as in \code{init:
can't open header for record REC}); \code{wfdberror()} returns the last such message, and
\code{wfdbquiet()} turns these messages off (\cref{sec:julia-library}).

\begin{table}[htbp]
\centering
\caption{The exceptions of the high-level API}
\label{tab:julia-exceptions}
\begin{tabularx}{\linewidth}{@{}lL@{}}
\toprule
\tabhead Exception & \tabhead Cases \\
\midrule
\code{WfdbError} & a record, header or annotation file cannot be found, read, created or
  written completely (for example on a full disk): ``can't read record REC'', ``can't read
  header for record REC'', ``can't write header for record REC'', ``write error on record
  REC'', ``directory D does not exist'', \ldots; a signal or annotation file ends before
  its header or its end-of-file marker says (``\ldots(getvec status S)'', ``\ldots(getann
  status S)''); an invalid argument of the readers and writers (most messages of
  \code{write\_record} and \code{write\_annotations} begin with the function's name,
  including those of an invalid \code{metadata} argument of \code{write\_record}); the
  operation needs a value that the record does not give: a time in seconds without a
  sampling frequency (``sampling frequency of record REC is not specified''), physical
  units without a gain (``signal N of record REC is uncalibrated'') \\
\code{ArgumentError} & an invalid argument of the metadata functions: an invalid key or
  value, a record name that is empty or \code{"-"}, a key given more than once in
  \code{changes}, a key that the header has more than once (to be set), a header without a
  record line, or a header read from a URL to be changed without \code{dir}; an empty
  record name or an annotator name that is not a string, in \code{check\_record} \\
\code{MethodError} & an argument of a type that no method accepts \\
\bottomrule
\end{tabularx}
\end{table}

\paragraph{Warnings} are logged with Julia's logging macro \code{@warn}, at the level
\code{Warn}, so the current logger decides where they go (the standard error by default);
\code{Test.@test\_logs} checks them, and a logger such as
\code{Base.CoreLogging.NullLogger()} silences them.

\paragraph{The checksum warning.} When \code{read\_record} reads to the end of a record
whose samples do not match the checksums of its header, it returns all the samples and
logs the warning ``checksum error in record REC: the signal data do not match the
checksums in its header''. \code{checksum\_warning=false} turns the warning off;
\code{checksum\_ok} of the result is \code{false} either way.

\paragraph{The duplicate-key warning.} When a header holds two metadata entries whose keys
differ only in letter case (written by other software), \code{read\_metadata},
\code{get\_metadata} and \code{read\_header} use the first and log the warning ``header of
record REC has more than one entry for key KEY (lines L1, L2); the first is used'' (KEY
spelled as in its first entry). \code{set\_metadata} and \code{update\_metadata} then
refuse to change that key (an \code{ArgumentError}), and \code{remove\_metadata} removes
all of its entries.

\paragraph{The library's settings after an error.} Whether a high-level function returns or
throws, the getvec mode (\code{getgvmode()}) and the time resolution of output annotators
(\code{getafreq()}) are those of before the call, as are legacy mode and the input and
output directories (\cref{sec:julia-conventions}). Opening an EDF record changes both, as
in C, and \code{wfdbquit()} resets neither: after \code{read\_header("100s.edf")}, a
\code{write\_annotations} without \code{fs} writes no time resolution, and the next
\code{read\_record} reads in the mode it reads in by default. The C-style functions
(\code{isigopen}, \code{wfdbquit}, \ldots) keep the behavior of C.

\begin{lstlisting}[language=Julia,style=example]
# A record whose last sample is changed after it was written: the checksum fails.
write_record("ck", reshape(Int32.(0:99), :, 1), 100; gains=100, adcres=12, units="mV")
dat = read("ck.dat"); dat[end-1:end] .= 0x00; write("ck.dat", dat)
rec = read_record("ck")         # logs the checksum warning
println(size(rec.samples, 1), " ", rec.checksum_ok)  # 100 false: all the samples
rec = read_record("ck"; checksum_warning=false)     # no warning

# Two entries for one key: the first is used, with a warning.
write_record("dup", zeros(Int32, 10, 1), 100; gains=0, adcres=12,
             comments=[" Age: 63", " AGE: 64"])
println(String(get_metadata("dup", "age")))         # 63 (with a warning)
try
    set_metadata("dup", "Age", "65")
catch e
    println(e.msg)
    # header of record dup has more than one entry for key Age (lines 3, 4); remove them first
end
println(remove_metadata("dup", "age"))              # 2

try
    read_record("nosuchrecord")     # (the library also prints a message)
catch e
    println(e isa WfdbError, " ", e.msg)            # true can't read record nosuchrecord
end
\end{lstlisting}

% ==========================================================================================
\section{The WFDB library functions in Julia}
\label{sec:julia-library}

Every function of the WFDB library is a function of the module \code{WFDBX}, under its C
name, with the same arguments, in the same order, and the same return values, so that
programs and documentation written for the C library carry over directly. The functions,
their arguments and their results are documented once, for all the languages, in
\cref{chap:library}; this section gives the rules by which the C declarations become
Julia functions (see also \cref{sec:library-julia}). The docstring of each function
(\code{?getvec}) summarizes it, and says what WFDB-x changes in it.

\subsection{Conventions}

\paragraph{The library state.} The state of the C library (the open record, its signals,
the annotators, the WFDB path, the time base, \ldots) is held by a \code{Wfdb} object, and
each C function that uses it takes one as an optional first argument:
\code{isigopen(w, "100s", si, 2)}, \code{getvec(w, v)}. Without it, the function uses the
default object, \code{lib()}. \code{Wfdb()} creates an independent one; separate objects
share no state, unlike the C library. The functions that use no state take none:
\code{wfdbldflags}, \code{wfdbcflags}, \code{wfdbdefwfdb}, \code{wfdbdefwfdbcal},
\code{isann}, \code{ammap}, \code{mamap} (and their \code{wfdb\_} forms), and
\code{wfdbx\_write\_errors} and \code{wfdbx\_note\_write\_error}, whose count belongs to
the process.

\paragraph{Names.} The functions, the constants (\code{WFDB\_READ},
\code{WFDB\_INVALID\_SAMPLE}, \code{WFDB\_HIGHRES}, \ldots) and the structures keep their C
names. \code{using WFDBX} makes most of them available; the others are called with the
module prefix: \code{WFDBX.wfdb\_error}, \code{WFDBX.wfdb\_me\_fatal}, the annotation codes
(\code{WFDBX.NORMAL}, \code{WFDBX.PVC}, \ldots, \cref{chap:formats}), and the function
forms of the macros of \file{<wfdb/ecgmap.h>} (\code{WFDBX.wfdb\_isann},
\code{WFDBX.wfdb\_isqrs}, \ldots; \code{isann}, \code{isqrs}, \code{map1}, \ldots{} are
exported).

\paragraph{Arguments.} String arguments are \code{String}s or bytes (\code{ByteStr}); a
byte string ends at its first NUL byte, as a C string does. C's \code{NULL} is
\code{nothing}: \code{sampfreq(nothing)} returns the frequency of the open record, and
\code{getinfo(nothing)} continues with the info strings of the last record. Where C fills
an array, the caller passes a vector: of \code{Siginfo} structures (\code{[Siginfo() for
\_ in 1:n]}), or of samples, which must be a \code{Vector\{Int32\}}; \code{nothing} stands
for C's \code{NULL} (\code{isigopen("100s", nothing, 0)}). Where C fills a structure, the
caller passes a mutable structure (\code{getann(0, a)} with \code{a = Annotation()}).
Where C fills an \code{int} or a pointer, the caller passes a \code{Ref}:
\code{getseginfo(r)} takes a \code{Ref}, which receives a vector of \code{Seginfo} (or
\code{nothing}).

\paragraph{Strings} are returned as bytes, \code{Vector\{UInt8\}}, or \code{nothing} where
the C function may return \code{NULL} (\code{anndesc}, \code{getinfo}, \code{wfdbfile},
\ldots), except the results of \code{timstr}, \code{mstimstr} and \code{datstr}, which are
\code{String}s; unlike C's pointers to static buffers, they stay valid after later calls.

\paragraph{Numbers.} The integer types are those of the C reference platform:
\code{WFDB\_Sample} and C's \code{int} are \code{Int32}, \code{WFDB\_Time},
\code{WFDB\_Date} and C's \code{long} are \code{Int64}, \code{WFDB\_Frequency} and
\code{WFDB\_Gain} are \code{Float64}, C's \code{char} is \code{Int8}, and
\code{WFDB\_Signal}, \code{WFDB\_Group} and \code{WFDB\_Annotator} are \code{UInt32}.
Julia's fixed-width arithmetic wraps around as C's does. Arrays are indexed from 1, but
the values that C counts from 0 (signal numbers, annotator numbers, sample numbers) keep
their C values.

\paragraph{Return values and errors.} The functions return what the C functions return:
an \code{Int32} status (negative on errors), a number of signals, a \code{Float64}, a
string or \code{nothing}. Errors are reported by these values and by a message on the
standard error, as in C, never by exceptions. \code{wfdberror()} returns the last message;
\code{wfdbquiet()} turns the messages off and \code{wfdbverbose()} on again.
\code{WFDBX.wfdb\_error(fmt, args...)} formats its message as C's \code{printf}.

\paragraph{State.} As in C, a \code{Wfdb} object keeps its state between calls: the open
record, its input signals, the output signals and the annotators remain open until
\code{wfdbquit()} closes them. \code{wfdbquit()} leaves the input and output directories
and the network setting as they are. The high-level functions of \cref{sec:julia-api}
(except the metadata functions) call \code{wfdbquit()} on the \code{Wfdb} object they use,
which closes anything opened with the library functions on that object.

\paragraph{Additions.} \code{getvecs(out, n)} reads up to \code{n} sample vectors at once
into \code{out}, a \code{Vector\{Int32\}} with room for \code{n} times the number of
values of one vector, with exactly the result of \code{n} calls of \code{getvec}, but
decoding whole blocks of the signal files; it returns the number of vectors read
($\ge 0$), or the negative status of \code{getvec} if the first call failed.
\code{getvecs\_status(out, n)} also gives the status of the last \code{getvec} call
(\cref{sec:library-additions}):
\begin{lstlisting}[language=Julia]
getvecs([w], out::Vector{Int32}, n) -> Int64
getvecs_status([w], out::Vector{Int32}, n) -> (Int64, Int32)
\end{lstlisting}
\code{WfdbFile} is C's \code{FILE *} (or \code{WFDB\_FILE}): \code{wfdb\_open} and
\code{wfdb\_fopen} return a \code{WfdbFile} or \code{nothing}, read and written with the
functions of the C standard library (\code{getc}, \code{fread!}, \code{fgets!},
\code{fwrite}, \code{fseek}, \code{fclose}, \ldots).

\subsection{The structures}
\label{sec:julia-structures}

The structures of the C library have the names of the C structures without the prefix
\code{WFDB\_}, and the C field names; they are mutable, with keyword constructors whose
defaults are C's zero values (\cref{tab:julia-structures}).

\begin{xltabular}{\linewidth}{@{}lL@{}}
\caption{The C structures in Julia}
\label{tab:julia-structures} \\
\toprule
\tabhead Structure & \tabhead Fields \\
\midrule
\endfirsthead
\toprule
\tabhead Structure & \tabhead Fields \\
\midrule
\endhead
\bottomrule
\endfoot
\code{Siginfo} & \code{fname}, \code{desc}, \code{units}
  (\code{Union\{Nothing,Vector\{UInt8\}\}}), \code{gain}, \code{initval}, \code{group},
  \code{fmt}, \code{spf}, \code{bsize}, \code{adcres}, \code{adczero}, \code{baseline},
  \code{nsamp}, \code{cksum} (\cref{tab:julia-siginfo}) \\
\code{Seginfo} & \code{recname} (the name of the segment, bytes), \code{nsamp} (its
  number of samples, \code{Int64}), \code{samp0} (the number of its first sample in the
  record, \code{Int64}) \\
\code{Anninfo} & \code{name} (the annotator name, bytes), \code{stat} (\code{Int32}:
  \code{WFDB\_READ} or \code{WFDB\_WRITE} for the MIT format, \code{WFDB\_AHA\_READ} or
  \code{WFDB\_AHA\_WRITE} for the AHA format); \code{Anninfo("atr", WFDB\_READ)} \\
\code{Annotation} & \code{time} (\code{Int64}, in sample intervals), \code{anntyp}
  (\code{Int8}), \code{subtyp} (\code{Int8}), \code{chan} (\code{UInt8}), \code{num}
  (\code{Int8}), \code{aux} (\code{nothing}, or bytes whose first byte is the length of
  the text that follows, as in C) \\
\code{Calinfo} & \code{low}, \code{high} (the levels of the calibration pulse, in physical
  units, \code{Float64}), \code{scale} (physical units per cm), \code{sigtype} and
  \code{units} (bytes or \code{nothing}), \code{caltype} (\code{Int32}) \\
\end{xltabular}

\subsection{The settings functions}
\label{sec:julia-settings}

The functions of \cref{tab:julia-settings} are additions of WFDB-x to the library: they
set where files are read and written, network access and legacy mode, for the
high-level functions, the library functions and the tools alike
(\cref{chap:files,chap:behavior}). Like the other library functions, each takes an
optional \code{Wfdb} object first, except \code{wfdbx\_write\_errors} and
\code{wfdbx\_note\_write\_error} (the count belongs to the process). Each \code{set}
function returns the previous setting, so that it can be restored. The C function
\code{setwfdb(p)} sets the input directory too (\code{p} must name at most one directory),
and \code{getwfdb()} returns it (\code{"."} if none is set); in legacy mode they set and
return the WFDB path, as in C.

\begin{xltabular}{\linewidth}{@{}lL@{}}
\caption{The settings functions (\code{[w]}: the optional \code{Wfdb} object)}
\label{tab:julia-settings} \\
\toprule
\tabhead Function & \tabhead Effect and result \\
\midrule
\endfirsthead
\toprule
\tabhead Function & \tabhead Effect and result \\
\midrule
\endhead
\bottomrule
\endfoot
\code{wfdbx\_set\_input\_dir([w], dir::OptByteStr)} & sets the input directory: a local
  directory, a \code{\textasciitilde} path or a URL; \code{nothing} or \code{""} restores
  the default (\code{WFDBX\_INPUT\_DIR}, or else the current directory). Returns the
  previous input directory as bytes, or \code{nothing} if none was set \\
\code{wfdbx\_get\_input\_dir([w])} & returns the input directory as bytes, or
  \code{nothing} if none is set \\
\code{wfdbx\_set\_output\_dir([w], dir::OptByteStr)} & sets the output directory: a local
  directory or a \code{\textasciitilde} path (a URL is refused, with a message, and the
  setting is unchanged); \code{nothing} or \code{""} restores the default
  (\code{WFDBX\_OUTPUT\_DIR}, or else the current directory). Returns the previous output
  directory, or \code{nothing} \\
\code{wfdbx\_get\_output\_dir([w])} & returns the output directory as bytes, or
  \code{nothing} if none is set \\
\code{wfdbx\_set\_network([w], on::Union\{Nothing,Bool,Integer\})} & turns network access
  on (\code{true} or $> 0$) or off (\code{false} or 0); \code{nothing} or a negative
  number restores the default (on if \code{WFDBX\_NETWORK} is a nonzero number). Returns
  the previous setting, 1 or 0 (\code{Int32}) \\
\code{wfdbx\_get\_network([w])} & returns 1 if network access is on, 0 otherwise \\
\code{wfdbx\_set\_legacy([w], on::Union\{Bool,Integer\})} & turns legacy mode on
  (\code{true} or nonzero) or off. Returns the previous setting, 1 or 0 \\
\code{wfdbx\_get\_legacy([w])} & returns 1 if legacy mode is on, 0 otherwise \\
\code{wfdbx\_legacy\_args([w], args::Vector\{String\})} & if \code{args[1]} (the
  arguments \emph{without} the program name, as Julia's \code{ARGS}) is
  \code{"-{}-legacy"}, removes it, turns legacy mode on and returns 1; otherwise returns
  0 \\
\code{wfdbx\_time\_needs\_freq([w], string::ByteStr)} & returns 1 if
  \code{strtim(string)} needs a sampling frequency, 0 otherwise \\
\code{wfdbx\_write\_errors()} & returns the number of write errors since the program
  started (output files that could not be written, flushed or closed); under
  \code{WFDBX.run\_tool}, since the tool started \\
\code{wfdbx\_note\_write\_error()} & counts a write error, for a program that writes a file
  itself \\
\code{wfdbx\_input\_path([w], name::OptByteStr)} & returns, as bytes, the path at which a
  program opens a file that it reads itself (not through the library): where
  \code{wfdb\_open} finds a file of that name (in normal mode, a relative name under the
  input directory, a \code{\textasciitilde} name expanded, an absolute name as given; in
  legacy mode, the name as given). \code{""} and \code{"-"} are returned as given;
  \code{nothing} for \code{nothing} or a URL (also one made by the input directory) \\
\code{wfdbx\_output\_path([w], name::OptByteStr)} & returns, as bytes, the path at which a
  program creates a file that it writes itself: where \code{wfdb\_open} creates a file of
  that name (in normal mode, a relative name under the output directory, a
  \code{\textasciitilde} name expanded, an absolute name as given; in legacy mode, any
  name under the output directory, used as a plain prefix). When the output directory has
  been prefixed, the missing directories of the path are created. \code{""} and
  \code{"-"} are returned as given; \code{nothing} for \code{nothing} or a URL \\
\end{xltabular}

\subsection{Examples}

Reading the samples and the annotations of a record, as a C program does (the record
\code{100s} is read from the current directory):

\begin{lstlisting}[language=Julia,style=example]
si = [Siginfo() for _ in 1:2]                       # receives one Siginfo per signal
nsig = isigopen("100s", si, 2)                      # open 2 signals
println(nsig, " ", String(si[1].desc), " ", si[1].fmt, " ", si[1].gain)  # 2 MLII 212 200.0
println(sampfreq(nothing))                          # 360.0: the frequency of the open record

v = zeros(Int32, 2)                                 # receives one sample per signal
isigsettime(strtim("0:30"))                         # skip to 30 seconds
getvec(v)
println(v[1], " ", v[2], " ", aduphys(0, v[1]))     # 947 968 -0.385

frames = zeros(Int32, 360 * 2)                      # the next second, in one call
println(getvecs_status(frames, 360))                # (360, 2)

annopen("100s", [Anninfo("atr", WFDB_READ)], 1)
a = Annotation()
n = 0
while getann(0, a) == 0                             # 0: an annotation was read; -1: the end
    global n += 1
    n <= 2 && println(mstimstr(a.time), " ", String(annstr(a.anntyp)))
end
#     0:00.050 +
#     0:00.214 N
println(n, " annotations")                          # 75 annotations

wfdbquiet()                                         # no messages on the standard error
stat = isigopen("nosuchrecord", nothing, 0)
print(stat, " ", String(wfdberror()))
# -1 init: can't open header for record nosuchrecord
wfdbverbose()
wfdbquit()                                          # close everything
\end{lstlisting}

Writing a record and an annotation file, sample by sample and annotation by annotation:

\begin{lstlisting}[language=Julia,style=example]
out = Siginfo(fname=b"ramp.dat", desc=b"ramp", units=b"mV", gain=200.0, fmt=16, adcres=12)
println(osigfopen([out], 1))                        # 1: one output signal
setsampfreq(250.0)
for k in 0:499
    putvec(Int32[k])
end
println(newheader("ramp"))                          # 0: ramp.hea is written
annopen("ramp", [Anninfo("qrs", WFDB_WRITE)], 1)
for t in (100, 350)
    putann(0, Annotation(time=t, anntyp=WFDBX.NORMAL))
end
wfdbquit()                                          # completes ramp.dat and ramp.qrs
println(length(read_annotations("ramp", "qrs")))    # 2
\end{lstlisting}

% ==========================================================================================
\section{The command-line tools in Julia}
\label{sec:julia-tools}

The 33 command-line tools (\code{rdsamp}, \code{rdann}, \code{wrsamp}, \code{xform},
\ldots) are functions of the package, one file per tool in \file{julia/src/apps}, run
through the driver script \file{julia/bin/wfdbx.jl}: its first argument names the tool,
and the other arguments go to the tool. Each follows its C program line by line. The
tools, their options, their output and their exit status are those of the C tools, and
are documented in \cref{chap:tools}. With the shell script \file{wfdbx} of
\cref{sec:julia-install}:

\begin{lstlisting}[style=shell]
export WFDBX_INPUT_DIR=/path/to/wfdb-x/tests/data
wfdbx rdsamp -r 100s -t 10 -p     # the first 10 seconds of 100s, in physical units
wfdbx rdann -r 100s -a atr        # its annotations
wfdbx rdsamp --legacy -r 100s     # legacy mode: --legacy as the first argument
wfdbx --version                   # WFDB-x 1.2.0
\end{lstlisting}

Without the script, the command is \cmd{julia -{}-startup-file=no
-{}-project=/path/to/wfdb-x/julia /path/to/wfdb-x/julia/bin/wfdbx.jl TOOL ARGUMENTS}.
Every run starts a new Julia process, which loads the precompiled package; most of the
start-up time of a tool is Julia's own. With no arguments, or an unknown tool name, the
driver prints a message on the standard error and exits with status 1. When an output
file or the standard output cannot be written (a full disk, for example), a tool reports
it (``PROGRAM: write error on standard output'') and exits with status 1; with
\opt{-{}-legacy}, it keeps the original exit status. A tool whose standard output is a
closed pipe ends with status 141, as a C program ended by \code{SIGPIPE}. The driver finds
the compiled package in the depot of the user's own home directory even when the
environment variable \code{HOME} is changed (\code{JULIA\_DEPOT\_PATH}, if set, is used as
given).

\paragraph{Inside a Julia session.} \code{WFDBX.run\_tool(tool, args; stdin=UInt8[],
dir=nothing, env=nothing)} runs a tool in the current process, with a fresh library state,
in the directory \code{dir}, with the environment variables of \code{env} set for the run
and the bytes \code{stdin} as its standard input, and returns its exit status and what it
wrote to the standard output and the standard error, as bytes. The caller's library state
and standard streams are restored when the tool returns. \code{WFDBX.tools()} returns the
names of the tools, and an unknown name is an \code{ArgumentError}.

\begin{lstlisting}[language=Julia,style=example]
status, out, err = WFDBX.run_tool("rdsamp", ["-r", "100s", "-t", "s3"]; dir=".")
println(status)                                     # 0
print(String(out))       # 3 lines: the sample number, then 995 and 1011 (tab-separated)
println(length(WFDBX.tools()))                      # 33
\end{lstlisting}

\paragraph{Environment variables.} The tools, like the rest of the library, use the
environment variables of \cref{tab:julia-environment}, also for the files that they read
or create themselves: the \opt{-i} file of \code{wrsamp}, \code{edf2mit},
\code{wav2mit}, \code{a2m} and \code{ad2m}, the input files of \code{ahaecg2mit} and the
\opt{-S} script of \code{xform} are read from the input directory (\code{wfdbdesc}
measures a signal file in the directory of its header), and the WAV file of
\code{mit2wav}, the EDF file of \code{mit2edf}, the master header of \code{snip -m} and
the \file{.des}, \file{.ecg}, \file{.key} and \file{.txt} files of \code{parsescp} are
created in the output directory, as \code{wfdbx\_input\_path} and
\code{wfdbx\_output\_path} locate them (\cref{sec:julia-settings}). Such a file cannot be
read from a URL. In legacy mode, the files that a tool reads itself are opened as named.

\begin{table}[htbp]
\centering
\caption{The environment variables used by the Julia library}
\label{tab:julia-environment}
\begin{tabularx}{\linewidth}{@{}lL@{}}
\toprule
\tabhead Variable & \tabhead Meaning \\
\midrule
\code{WFDBX\_INPUT\_DIR} & the input directory (a directory, a \code{\textasciitilde}
  path or a URL), when \code{wfdbx\_set\_input\_dir} does not set one \\
\code{WFDBX\_OUTPUT\_DIR} & the output directory, when \code{wfdbx\_set\_output\_dir} does
  not set one \\
\code{WFDBX\_NETWORK} & network access is on if it is a nonzero number, when
  \code{wfdbx\_set\_network} does not set it \\
\code{WFDBANNSORT} & \code{0} to keep annotation files in the order written (no sorting
  when they are closed) \\
\code{WFDBGVMODE} & the default getvec mode (\code{setgvmode} with a negative mode) \\
\code{WFDBCAL} & the calibration file read by \code{calopen(nothing)} (default
  \file{wfdbcal}) \\
\code{WFDB} & the WFDB path, used in legacy mode only (outside legacy mode, it is
  ignored, with a message) \\
\code{HOME} (\code{USERPROFILE} on Windows) & the home directory, for
  \code{\textasciitilde} paths \\
\bottomrule
\end{tabularx}
\end{table}

Tools that write signal parameters require them as options: \code{wrsamp} (with
\opt{-o}) requires \opt{-F} (frequency), \opt{-G} (gain), \opt{-R} (ADC resolution) and,
for calibrated signals, \opt{-U} (units); \code{wav2mit} requires \opt{-G}, and \opt{-U}
unless the gain is 0; \code{ad2m} and \code{ahaecg2mit} require \opt{-F}, \opt{-G},
\opt{-R}, and \opt{-U} unless the gain is 0; \code{a2m}, \code{rdann}, \code{ann2rr},
\code{sortann} and \code{m2a} require \opt{-F} when the record's header gives no
frequency (\cref{chap:tools}).

% ==========================================================================================
\section{Testing}
\label{sec:julia-testing}

The unit tests of the Julia implementation are in \file{julia/test}, and use the standard
library \code{Test}. From the root of the repository:

\begin{lstlisting}[style=shell]
julia --startup-file=no --project=julia julia/test/runtests.jl   # all the unit tests
julia --project=julia -e 'using Pkg; Pkg.test()'                  # the same, through Pkg
\end{lstlisting}

The tests cover the C standard functions (\code{printf}, \code{strtol}, \code{strtod},
\code{sscanf}, \ldots, against glibc), the files and the library functions
(\file{runtests.jl}, \file{signals.jl}, \file{test\_annotations.jl}), the FLAC codec
(\file{test\_flac.jl}), finding files, the settings and network access
(\file{test\_files.jl}), the high-level functions, the metadata functions
(\file{test\_metadata.jl}, with the cases shared by the seven languages, in
\file{tests/metadata}), \code{check\_record} (\file{test\_check.jl}) and the tools run
in-process. They run in temporary copies of the test records of \file{tests/data}. The
tests of network access use a local HTTP server on 127.0.0.1 only (written with the
standard library \code{Sockets}). The tests of \code{check\_record} write the corpus of
\file{python/tests/check\_corpus.py} with \cmd{python3}, NumPy and the Python package of
\file{python/src}, and compare the reports with its \file{expected.txt}; they are skipped
if that script cannot run.

The cross-language conformance suite runs every tool and compares its results, byte for
byte, with those of the original C programs (1135 cases), and, with \opt{-{}-legacy}, in
legacy mode (883 cases). It runs the tools through \file{julia/bin/wfdbx.jl}:

\begin{lstlisting}[style=shell]
PYTHONPATH=$PWD/python/src python3 tests/conformance/harness.py check --impl julia -j8
PYTHONPATH=$PWD/python/src python3 tests/conformance/harness.py check --impl julia -j8 --legacy
\end{lstlisting}
